
\typesetchapter{Chapter 13}{Horizon}

\textit{The Mountain will feel purer from here. Do not believe that.}\par
\textit{Distance flatters prisons and homelands equally.}\par
\textit{--- Ka-Sarah}\par

\ornament

The escape is what a listener wants first: the desert, the metal bird, the open water beyond a window. I will give all of it, in order, and slower than the wanting would like. But movement was the smallest part of leaving.\par

First they erased me.\par

\ornament

I stood in the sealed route beneath the hearing chamber while Justice removed me from the ordinary sequence of the Mountain. The public form of it passed without me, if public was the right word for a statement spoken before five witnesses in a room no family would enter. I heard only the lower version, given by the white-haired Justice witness while Shield prepared the transfer.\par

``Va-Kailan of Light ceases ordinary standing.''\par

The witness kept her eyes on the tablet as she read. That was mercy, or perhaps only discipline.\par

``Office consequence passes to Light under Keeper management. Household consequence passes to Faith under sealed wording. Material consequence passes to Balance. Movement consequence passes to Shield. Sequence consequence remains with Justice. The name is held from common use.''\par

Held.\par

Alive. Erased. Held.\par

A name in the Mountain was ration, bed, task, prayer line, training record, family place, work expectation, burial ledger. To hold a name was to lift it out of the channels through which a person remained among the living.\par

I looked at the dark cord Va-Sheva had tied below the covered white knot on my wrist. My Veil mark remained beneath Shield cloth. The Mountain made what I was unusable.\par

``Will my mother be told I am dead?'' I asked.\par

The witness turned one line on the tablet.\par

``Faith will give the household a form sufficient for survival.''\par

``That is not an answer.''\par

``It is the lawful answer.''\par

Va-Sheva stood three paces away, silent. No weapon showed on her. No explanation either.\par

``What form?'' I asked.\par

The Justice witness looked at me then.\par

``One they can repeat without committing breach.''\par

I regretted the question. A form they could repeat. A shape designed for mourning without knowledge. My mother would receive words she could survive. My father would understand the absence of details better than most men and be ruined by that understanding. My household would wonder whether to grieve death, disgrace, selection, or disappearance. Faith would stand with them. Faith would keep the wound from speaking too much.\par

That was mercy too, perhaps. I was beginning to hate how often mercy arrived wearing the face of structure.\par

The transfer passed no public gate.\par

Shield took me through routes that had no childhood memory attached to them. Below the hearing chamber, below the old water lines, below even the places where workers joked that the Mountain had no further patience for stairs, the passages became narrower and less finished. Some were cut by tool. Some had been enlarged from natural seams. Some were so low that even Va-Sheva had to bend. The air moved strangely there: thin, dry streams that came from cracks the official maps of ordinary people omitted.\par

Balance had prepared a transfer pack. It contained nothing sentimental. Dried fruit. Salted strips. Two water skins. A folded robe of rougher cloth. A mouth covering. Sand wraps. A small packet of bitter paste for the gums against thirst. A list of measures tied to the pack with green thread. Even exile had weight, water, salt, and inventory.\par

I noticed a second pack on one warden's back.\par

``Is that mine too?''\par

Va-Sheva answered without turning. ``No. That is what Shield carries because Balance assumes you may become useless.''\par

I almost smiled. It sounded like her, and few things left still did.\par

Justice stopped us at three points. Each time the witness marked my body, route, transfer condition, and silence. No one asked whether I wished to continue. Exception had already answered that question by removing its relevance.\par

At the last interior threshold, Faith waited. A single attendant of middle age with a covered face and a bowl of water in both hands. She was a stranger to me, or had been chosen because she could appear so. The bowl was shallow. The water inside held still.\par

She spoke no scripture. That omission hurt. Instead she said, ``What leaves under seal leaves under God.''\par

I lowered my head. ``What returns under seal returns only if God permits return.''\par

I looked up. The attendant touched water to my forehead, then to the covered knot at my wrist, then to the dark cord Shield had tied beneath it.\par

``What is held from the people is not hidden from God.''\par

Her hands withdrew. The blessing ended there. Short, almost bare, shaped so that no forbidden fact entered holy language. Faith could mourn without knowing. Faith could bless without being told what it blessed.\par

The final Shield door opened outward.\par

Night waited.\par

\ornament

He has taken four days to walk me down into the Mountain, and now, in a handful of quiet sentences, the Mountain has set him outside itself and held his name behind him like a door. He stops there, on the threshold, the way a man stops at the top of stairs he has no wish to descend.\par

``They have finished with you,'' I say. ``Good. Then give me the road.'' I keep it plain, an old man asking for the last of a long story before the fire dies. ``The way out --- all of it, in order. The desert, the well, the water. I have kept you from your bed four nights for the account; I would not stop it one door short of where it ends.''\par

``It ends badly,'' he says.\par

``Most roads out of a home do.'' The reels turn between us. I know where this road ends before he walks it, because it ends where I am sitting --- in this country, this house, at this table, on my side of it. He is about to describe reaching Horizon, and Horizon is the name across my own contract; there is a strangeness in waiting for a man to cross a desert to arrive at the very people who pay me to listen to him do it. He believes he is telling me how he lost the Mountain. I think he is telling me how he found people who mean to keep it --- and only a small, new voice in me, four days old, asks quietly whether those are the same sentence.\par

``Then walk,'' I say, ``and do not spare the turns.''\par

\ornament

The desert entered first as cold. Then height. Then smell: mineral, dry scrub, distant dust, something faintly animal beyond the prepared route. I stepped out beneath the Mountain's shadow for the second time in my life and the first time without prayer access. The sky above me was full of stars. No vault. No ceiling. No cut stone receiving sound. Nothing to teach the eye where to stop.\par

I hated it less this time. That frightened me.\par

The party moved before dawn. Va-Sheva led until the first ridge, then gave command to a Shield warden whose face had been wrapped so fully that only his eyes remained. Behind us came two more wardens and the Justice witness. No Faith. No Balance. Their parts had been completed at the threshold.\par

For a while the Mountain remained behind us, blacker than night. Then the path folded, and it vanished. I stopped. A warden touched my shoulder. Lightly. Enough. I turned back, but there was nothing to see. The Mountain had kept its place. The world had simply arranged itself so that the only home I had ever known no longer existed from where I stood.\par

``Move,'' Va-Sheva said. I moved.\par

At first the journey was silence, cold, and pain in the feet. The sand wraps helped a little. My body was trained for stone, stairs, chambers, controlled routes, measured heat. The desert had no interest in Mountain competence. It turned ankles, stole direction, widened distance, and made every ridge into a false promise.\par

When dawn came, Shield put us under stone. A cut in the rock barely deep enough for bodies and packs. A cloth of desert colour was drawn across the outer edge and pinned with stones. From outside, we would become shadow. Inside, we became breath, heat, thirst, and waiting.\par

I was given two mouthfuls of water. I wanted more and kept silent.\par

Va-Sheva sat near the opening, looking outward through a gap in the cloth. She had kept her gloves on. Dust had settled along the line of her cheek. Outside light found her face and made it older.\par

``Is this still the Mountain?'' I asked.\par

``No.''\par

``Is it still Shield?''\par

``Yes.''\par

That distinction was almost the whole of my new life.\par

The day passed in fragments. Heat rose. The cloth breathed outward and inward with small changes in air. Once, a warden moved to adjust a stone and froze for three full breaths before finishing. Something had passed beyond the cut, unseen by me. Animal, bird, man, wind. Shield narrated no danger for the comfort of those protected from it.\par

Near evening, Va-Sheva gave me a strip of salted meat. I took it.\par

``Did you know,'' I asked, ``before I showed you the scrap?''\par

She kept her eyes on the gap. ``You are asking which part.''\par

``Yes.''\par

``I knew there were outer holdings.''\par

``Holdings.''\par

``Yes.''\par

``Archives?''\par

She said nothing.\par

``Contacts?''\par

Still nothing.\par

``Sarah Blackwood?''\par

Her face changed so little I almost missed it. ``You may not say that name again on this route.''\par

``But you know it.''\par

``I said you may not say it again.''\par

The old anger rose in me, but it found no easy shape. The trial had burned away the simplicity of accusation. She knew something. Too little. Or enough but under seal. Perhaps knowledge in Shield came as passages did: by station. One warden knew the outer marker. One knew the outcrop. One knew the route below it. One knew a name without a face. One knew a face without a name. One Veil might know there were holdings and still lack standing to name Horizon.\par

Let no one hand hold the whole. Even guilt was divided.\par

``I thought knowing would make things clearer,'' I said.\par

``That is because Light lies to its children differently from Shield.''\par

I looked at her. ``How does Shield lie?''\par

``We tell ours that if they train enough, they will see danger before it harms what they love.''\par

``And that is false?''\par

She turned from the cloth gap then. ``It is true often enough to destroy them.''\par

The answer silenced me.\par

At dusk we moved again.\par

The second night brought us farther from the Mountain than I knew how to measure. The stars wheeled overhead. Shield used them without naming them to me. The path avoided obvious valleys and easier ground, choosing instead hard shelves, dry cuts, stone ridges where footprints broke or vanished. Once we crossed a field of pale rock where every step seemed too loud. Once we lay flat while a distant engine passed beyond a ridge, a sound so alien that I first thought it some desert animal too large for scripture.\par

``What is it?'' I whispered.\par

Va-Sheva's hand closed once in warning. I obeyed. The engine faded. Only later, when we moved again, she said, ``Vehicle.''\par

The word belonged to no chamber I knew. ``Outsider?''\par

``Yes.''\par

``Did it see us?''\par

``No.''\par

I believed her because I had to. That dependency angered me more than ignorance.\par

Before the third dawn, the route ended at a dry well. It looked abandoned. Stones fallen, lip cracked, rope long gone. A dead thorn tree leaned over it like a warning no one maintained. Shield circled the place before approaching. Va-Sheva inspected the inner wall with two fingers, found something, and pressed.\par

From below came a hollow answer. A message travelling through old depth. We waited.\par

After a long while, a voice from inside the well spoke in a language strange to me. One of the wardens answered in the same language. The stone at the base shifted.\par

Horizon began as a hole beneath a dead well. That disappointed me until I understood that disappointment was the child's expectation of grandeur surviving where the man had no use for it.\par

The passage below the well was wider than Shield's desert route, cooler, and prepared for repeated movement. Its walls were reinforced with outsider materials: metal ribs, smooth white panels, rubber-smelling seals, small glass bulbs that gave steady light without flame. I stopped before the first one despite myself.\par

``No lamp,'' I said.\par

A man standing inside the passage laughed softly. ``Not in your sense.''\par

The accent was strange. The language was Mountain speech, but bent differently, as if each word had crossed water and returned changed. The man was tall, broad-shouldered, and perhaps in his early forties. He wore desert-coloured clothes unlike Shield leather, fitted for movement rather than ceremony. His hair was dark with grey at the sides. His face was lined by sun and discipline. A short blade sat at his belt beside stranger tools of metal and black material. He looked at the wardens, then at Va-Sheva, then at me.\par

``Jean-Charles,'' Va-Sheva said.\par

The man inclined his head. ``Va-Sheva.''\par

There was history in the way he said her name. Recognition with more warmth than duty required. I saw it before I wished to see it, and because I saw it, something mean and young stirred in me even there, even then, beneath a dead well on the first route of exile. I despised myself for it.\par

Jean-Charles looked at me and switched into more careful Mountain speech. ``So this is him.''\par

``This is the transfer,'' Va-Sheva said.\par

``Of course.'' The correction was mild but received. He stepped aside. ``Welcome to the part of the world that pretends it is not also a cave.''\par

I said nothing.\par

Jean-Charles smiled as if silence amused rather than offended him. ``Good. Still alive inside the head, then.''\par

Va-Sheva said, ``He is under Exception.''\par

``I know.''\par

``Do not test him for sport.''\par

``I never test for sport.''\par

The lie was so effortless that even the wardens seemed to ignore it by long practice.\par

We moved deeper.\par

At the first inner chamber, the Justice witness handed a sealed packet to a woman waiting beside a metal table. The woman had pale hair cut shorter than any Mountain woman's, and eyes that assessed without asking permission. She wore the grey of outsider cloth rather than of Candidates: jacket, trousers, boots, a narrow cord at the throat bearing no Kahir sign. Her face was neither young nor old in Mountain terms. She had the composure of someone who belonged to another hierarchy entirely.\par

``Elise,'' Jean-Charles said.\par

The woman ignored him and took the packet.\par

Va-Sheva addressed her formally. ``Va-Elise.''\par

So the title remained here too, though its meaning had changed by crossing.\par

Va-Elise broke the outer seal under the witness's eyes, read the first page, left the second, and signed a mark in ink on a smooth white sheet. Ink from a tube. Paper too perfect. The objects themselves were small shocks.\par

``Body received under Shield-Horizon transfer,'' she said. ``Speech restricted. Return unpresumed. Name held. Public standing ceased.''\par

The Justice witness answered, ``Sequence delivered.''\par

Va-Elise looked at me for the first time. ``You understand your old name is not for ordinary use here.''\par

``No.''\par

``Good. Then listen. Va-Kailan belongs to the Mountain record. Kailan belongs to what you were before rank. Outside tongues will damage it immediately. For operational use, you will be Kylian. Within Horizon classification, because you are under instruction, not command, you are Sa-Kylian.''\par

The demotion entered more deeply than I expected. ``Candidate?''\par

``Provisional Candidate.''\par

``I was Veil of Light.''\par

Va-Elise's expression held. ``In the Mountain. Under a title now sealed from common use. Here you do not know how a door opens, how money works, how a public street is crossed, how a lie is printed, how a telephone betrays location, how a camera steals sight, how a church arranges memory, how a passport creates a person, or how an Outsider smiles while measuring whether you are useful. Candidate is generous.''\par

Jean-Charles murmured, ``She likes to welcome people warmly.''\par

Va-Elise let that pass.\par

``She should not be contradicted,'' Va-Sheva said to me. The sentence landed strangely: Va-Sheva was instructing survival.\par

I --- no, Sa-Kylian, if the new grammar had already begun --- looked at the table. There were objects arranged on it: folded clothes, dark shoes, a small book waiting for my image, metal keys, a plastic card, a map printed in colours too precise to be a gravure. Beside them lay a pair of spectacles with clear lenses.\par

``I do not need those,'' I said.\par

Va-Elise followed my gaze. ``They are not for correction. They are for concealment.''\par

``Of what?''\par

``Attention. People look differently at men with spectacles.''\par

That was absurd. I looked at Jean-Charles. The man nodded gravely. ``Sadly true.''\par

The outer world was going to be worse than death in ways no scripture had prepared me for.\par

Va-Sheva removed the dark cord from my wrist. The cord only. She gave it to Va-Elise, who placed it inside a small envelope and marked the time. The gesture made the transfer physical.\par

``Am I still under Shield?'' I asked.\par

Va-Elise answered before Va-Sheva could. ``You are under Horizon custody by Shield transfer and Justice sequence.''\par

``What is Horizon?''\par

``A line you cannot reach by walking toward it.''\par

Jean-Charles smiled.\par

Va-Elise continued, ``That is the poetic answer. The practical answer is: the structure that keeps certain Mountain facts from dying inside the Mountain and certain outside facts from entering the Mountain uncontained.''\par

``Another Kahir?''\par

``No.''\par

Va-Sheva said, ``Do not try to fit it into five hands.''\par

That frightened me more than the dead well. The five hands had explained the world with a shape my mind could inhabit. If Horizon escaped the five, then Horizon was either outside Covenant or built in a gap Covenant had prepared and hidden.\par

``Who commands it?''\par

Va-Elise closed the packet. ``For now, Ka-Sarah.''\par

Sarah. The name returned before I could prevent it. Sarah Blackwood under desert light. The photograph. The record. The word Horizon sealed near her name. I said nothing.\par

Va-Elise saw enough. ``You will meet her when you are clean, dressed, fed, and no longer looking as if every object in the room is a theological insult.''\par

``They are.''\par

Jean-Charles laughed. Va-Elise's mouth moved by a fraction. ``Good. He is not stupid.''\par

Va-Sheva said, ``No.''\par

There was something in that no. Certainty. It reached me before I could defend against it.\par

\ornament

They gave me water, food, and clothes.\par

The clothes were the first humiliation, and ugliness was the least of it. Trousers, shirt, jacket, underclothes shaped without modest ritual, socks, hard-soled shoes. They made me feel both exposed and disguised. Mountain robes arranged the body by office. These clothes arranged it by usefulness in a world of strangers.\par

Va-Sheva stayed away while I changed. That was mercy.\par

When I returned, Jean-Charles looked me over and said, ``Almost human.''\par

I stared.\par

``In the outside sense,'' he added.\par

``That does not improve it.''\par

``No. But you understood the insult. Progress.''\par

Va-Elise handed me the spectacles. I put them on. The clear lenses caught the white electric bulbs and placed small false lights at the edge of my sight. ``I hate them,'' I said.\par

``They suit you,'' Jean-Charles said.\par

Va-Sheva entered then, saw me, and stopped. For the first time since the hearing, she nearly smiled. Too little for anyone else to accuse her of it. Enough for me to be wounded by the kindness.\par

``You look like an Outsider scholar,'' she said.\par

``I have been called worse today.''\par

``Yes.''\par

The word changed the room again. The transfer had reached its final personal moment and everyone there knew it. Va-Elise became suddenly occupied with the packet. Jean-Charles inspected a tool that required no inspection. The remaining warden looked at the wall.\par

Va-Sheva came closer. ``This is where Shield leaves you.''\par

``You leave me.''\par

``Yes.''\par

``Do you return to the Mountain?''\par

``Yes.''\par

``Will you tell them I reached Horizon?''\par

``No.''\par

``Will Ka-Raedin know?''\par

``Through sequence. Not through me.''\par

``And my family?''\par

She said nothing. I looked away first. That was new. She had often looked away first when we were younger, because she had somewhere to go and I had nowhere she permitted me to follow. This time I looked away to keep from asking for something neither of us could lawfully give.\par

``You did not choose death,'' I said.\par

``No.''\par

``You chose this.''\par

``I chose the only containment that left you breathing.''\par

``That is not the same as choosing me.'' The sentence escaped before discipline could stop it.\par

Va-Sheva received it without visible movement. Jean-Charles became very still across the room. Va-Elise kept her eyes down, which meant she was listening completely.\par

Va-Sheva answered quietly. ``No.''\par

It was cruel because it was exact.\par

I nodded once. ``Then thank you for choosing breathing.''\par

Her face changed then. Pain, almost. Or anger. Or something too divided by office to survive in the open. ``You are alive,'' she said.\par

``So I keep being told.''\par

``Believe it slowly.'' Her hand moved toward the covered knot at my wrist, then stopped before touch could become promise. She stepped back before the moment became something Shield would have to contain. Then she left.\par

No embrace. No private token. No promise. The door closed behind her, and with it the last person from the Mountain who had known me as Kailan before titles began cutting pieces from the name.\par

I stood in outsider clothes beneath electric light and understood that survival had brought me beyond the language in which punishment could be recognised.\par

\ornament

He has reached us. In his own telling he is standing in a room beneath a dead well with the dust of the desert still on him, and the names he lays on the table are ours --- Jean-Charles at the lamp, Elise with her packet and her cold ledger, and above them, in a house above a lake still ahead of him in the account, the name he lowers his voice to say. I have known all three of them longer than he has been alive to the outside, and it is the strangest hour of four strange days to hear them handed back to me young --- met as fate, as figures in a hidden people's story, rather than as the colleagues who share my corridors.\par

``You know them,'' he says. The account has stopped; he has lifted his eyes from the reels and set them on me. ``You told me what you are. So you know these names.''\par

``I know them.'' It costs me nothing to say so, and after four days I have stopped hoarding the things that cost me nothing. ``Jean-Charles is exactly as tiresome as you found him. Elise has not warmed once in all the time I have known her.'' He almost smiles, and the almost is worth more to me than I expected any part of this work to be worth.\par

``And her,'' he says. ``Ka-Sarah.''\par

``The account first.'' I say it gently --- only a fact-checker's habit: I would rather have the road whole than let either of us leap ahead to the woman at the end of it. But for the length of a breath he holds my eyes, and I let him, and I wonder, tonight, which of us sits nearer his own doubt --- he to his, or I to mine. ``Finish the road,'' I say. ``Bring yourself into the house above the water. I want the first time you saw her.''\par

\ornament

Instruction came before Ka-Sarah. A room rather than a chamber, filled with objects more dangerous for appearing harmless. A telephone. A camera. A printed newspaper. A map of the world, obscene in size. Coins and paper money. A passport with a blank place where my image would be placed. Va-Elise taught quickly and without pity. Jean-Charles interrupted only to translate the practical stupidity of things.\par

``A car is a small room that moves faster than a horse,'' he said.\par

``I know what a vehicle is now.''\par

``No. You know the word. You do not yet know that roads are rivers made of decisions.''\par

Va-Elise said, ``That one is actually useful. Remember it.''\par

They moved me again before dawn. This time by vehicle. The first engine vibration made me grip the seat hard enough for Jean-Charles to notice. ``You are safer than you feel.''\par

``That is not a meaningful category.''\par

``Good answer.''\par

We crossed dark land, then roads, then places where outside lights multiplied without order. Towns. Buildings. Wires. Signs. Machines. People moving without visible permission. A world mad with doors and no wardens at them. I said little.\par

The journey after that became a sequence of enclosed transitions: vehicle, hidden room, aircraft, vehicle again, another room. The aircraft was the worst. To be above the earth, between Heaven and ground, held inside a metal body full of strangers who ate, slept, coughed, read, complained, and trusted height as if it were lawful --- that broke something ordinary in me. Jean-Charles sat beside me and said only, ``Do not look down until I say.'' For once I obeyed without resentment.\par

We came to Switzerland in the spring of 1998.\par

I learned the date from a newspaper before I learned how cold the air would feel. And here the telling must slow, for this is the country in which it is now being told --- a few valleys from the room where I sit tonight. Everything before this was the story of a boy leaving a cave. This is the part that has never entirely finished happening.\par

Switzerland kept nothing of the desert. That was the first betrayal of expectation. The outside was many things at once: mountains without secrecy, water without Covenant, roads cut openly through stone, houses exposed to sky as if sky were harmless, churches with towers that pierced upward instead of hiding below. Snow remained on distant peaks though the lower ground had begun to soften. The air smelled of water, metal, pine, fuel, bread, and strangers.\par

They placed me in a house above a lake. A house. That was its cruelty. The windows opened. The doors had locks I could understand and others beyond me. Food came in packages printed with languages. Water came by turning a metal handle. Light came by touching a switch. No one prayed when it appeared. No one measured the first cup. No one thanked the Mountain because there was no Mountain to thank.\par

For three days I barely slept. On the fourth, Ka-Sarah received me.\par

She was older than the photograph and still recognisable enough that recognition became a form of vertigo. Sarah Blackwood had been young under desert light; Ka-Sarah was old. Her hair had gone mostly white, gathered at the back of her head. Her face carried age without surrendering command. She wore dark wool instead of robes, and a small gold cross at her throat so discreet that it seemed more dangerous than display. Her eyes were blue-grey, cold, amused, and very tired.\par

The room where she waited had books on three walls, a map of North Africa behind glass, and a window overlooking the lake. No curtains were drawn. She sat with her back partly to the open view, as if no sane person would consider light from outside a tactical error.\par

Va-Elise stood beside the door. Jean-Charles near the window. No Va-Sheva.\par

I stood. The correct movement escaped me.\par

Ka-Sarah observed me with open interest. ``It is not every decade that Horizon receives a Candidate,'' she said. ``An ex-Veil, no less. How extravagant of Shield.''\par

The line was elegant enough to be insulting and amused enough to be worse. ``I was Veil of Light,'' I said.\par

``Yes,'' she said. ``That is what ex means.''\par

I disliked her at once. That seemed to please her.\par

``You are Sarah Blackwood.''\par

``I was Sarah Blackwood before most of the men who feared me learned how to spell it. Here, you may call me Ka-Sarah.''\par

``Keeper of Horizon?''\par

``Among other inconveniences.''\par

``Horizon is not a Kahir.''\par

``No. But Kahir grammar is hard to remove from children raised in caves.''\par

Va-Elise said nothing. Jean-Charles looked out the window with the expression of a man enjoying weather.\par

Ka-Sarah gestured to the chair before her. ``Sit. You are making exile look theatrical.''\par

I sat. She studied me over folded hands.\par

``You have seen a room you were not meant to see, taken proof you were not meant to carry, survived a trial that had excellent reasons to kill you, crossed the desert, flown over Europe without screaming, and now you are in Switzerland wearing spectacles intended to make you look harmless. An eventful week.''\par

``Why am I here?''\par

``Because dead boys tell us less.''\par

``I am not a boy.''\par

``No. That appears to be part of the difficulty.''\par

She opened a file on the table. Paper, photographs, typed sheets, copied Mountain marks. My life had already become outside paperwork.\par

``You are under Exception. That means you are alive by procedure, not by freedom. Do not confuse the two. You will not leave this property without permission. You will not speak of the Mountain to anyone outside authorised Horizon structure. You will not attempt to contact your family, Ka-Raedin, Va-Sheva, or any Mountain office. You will not return toward the desert. You will not handle maps without supervision. You will not use telephones. You will not keep private writing materials until we decide whether your mind can hold paper responsibly.''\par

I almost laughed at that. The sound stayed down.\par

``And,'' Ka-Sarah said, ``you will not enter churches.''\par

That sentence landed differently. ``Why?''\par

Her eyes sharpened. ``Because you are too full of broken scripture, and churches are machines for arranging memory in public. You are not ready to stand inside one without either worshipping, condemning, or stealing from it. All three would be inconvenient.''\par

I thought of the tortured man on crossed wood. The father's blade. The many-handed mother. The death-headed figure. The false water lord. I thought of Faith beneath stone and Ka-Elun's voice reading Covenant. I thought of Va-Maera saying that a valve disproved nothing of God.\par

``You wear a cross,'' I said.\par

Ka-Sarah touched it lightly. ``Yes.''\par

``Are you one of them?''\par

``One of whom?''\par

``Outsiders.''\par

She smiled then. ``My dear boy, from where you sit, everyone is.''\par

I had no answer. She closed the file.\par

``You think you have been taken beyond the Mountain. You have not. You have been moved to the Mountain's horizon --- the place where its shadow meets the world it refuses to name properly. You are not free. You are not dead. You are not restored. You are not yet useful. That is your present condition.''\par

``And if I refuse it?''\par

``Then Shield was right and Justice was merciful for too long.''\par

There was no cruelty in her voice. Only administrative fact.\par

I looked past her, through the window, toward the lake. Open water under an open sky. Boats crossing in full view; houses along the shore, each with its windows, its lights, its doors, its name. No shadow of Mountain. No Oasis hidden from Heaven. No threshold drawn between safety and the curse.\par

Another man would have called it freedom. All I could feel was how much of me it left uncovered.\par

``What do you want from me?'' I asked.\par

``For now? To survive the outside without turning every object into theology.''\par

``And after that?''\par

Ka-Sarah leaned back. ``After that, perhaps you will tell us what Light saw before Shield buried you.''\par

``You already know.''\par

``No,'' she said. ``We know what Shield could afford to transfer. We know what Justice could lawfully write. We know what Light is probably hiding from itself. We do not know what you saw.''\par

That sentence changed the room. For the first time since the trial, my sight became something other than offence. I distrusted the relief immediately.\par

Ka-Sarah saw that too. ``Good,'' she said. ``You are learning.''\par

I almost hated her for using the word in Ka-Raedin's rhythm.\par

She stood, and the interview was over because she had decided it was. ``One more thing, Sa-Kylian.'' I looked at her. ``The Mountain will feel purer from here. Do not believe that. Distance flatters prisons and homelands equally.''\par

Then she left me with Va-Elise and Jean-Charles, in a room full of books forbidden me, before a lake forbidden me alone, under a sky no one had measured for me.\par

No one spoke for a while. Then Jean-Charles said, ``Well. That went beautifully.''\par

I looked at him. ``It did not.''\par

``No,'' Jean-Charles said. ``But she did not have anyone shot, which in this house counts as hospitality.''\par

Va-Elise said, ``Do not make him think that is a joke.''\par

Jean-Charles considered this. ``Half a joke.''\par

The absurdity reached me late. As a small failure of misery to keep perfect form. Va-Elise noticed and let it stand. That was either kindness or efficiency. With Horizon, I had already learned to leave the two joined.\par

They fed me after that. Food. Bread with a crust that broke under the hand. Cheese sharp enough to make my jaw tighten. Soup with herbs new to me. An apple, cut into quarters because Va-Elise had decided I had yet to earn the small knife placed beside the plate. Water in a glass so clear it looked like a trick of Light.\par

I waited for the prayer. None came.\par

Jean-Charles lifted his own glass. ``Food counted.''\par

Va-Elise looked at him. ``Do not.''\par

``Witnessed,'' Jean-Charles added.\par

Discipline forbade the laugh. I laughed. Once. Quietly. It hurt my throat. The laugh freed nothing. It forgave no one. It stopped short of the Mountain, my mother, Ka-Raedin, and Va-Sheva crossing the desert back into stone. It merely proved that the body, treacherous and practical, sometimes accepts warmth before the mind has granted permission.\par

After the meal, Va-Elise placed a narrow book on the table before me. Its cover was plain grey. The title had been written twice: once in an outside language beyond my reading, and once in careful Mountain hand.\par

\textit{On Light, Lenses, and the Conduct of Images.}\par

I left it where it lay. ``What is this?''\par

``An optics manual,'' Va-Elise said.\par

``I can see that.''\par

``Then why ask?''\par

``Because I do not understand why it is being given to me.''\par

Jean-Charles leaned back in his chair. ``Because Horizon is cruel in more imaginative ways than Shield.''\par

Va-Elise ignored him. ``It has been translated into Nocens where translation was possible. Where it was not possible, we left the outside term and added notes.''\par

I looked at the book as if it might accuse me. ``Am I permitted to read it?''\par

``Under supervision.''\par

``Of course.''\par

``You were Veil of Light,'' Va-Elise said. ``That is precisely why we do not intend to leave you alone with a book that explains what Light has not yet named.''\par

That should have offended me. It did offend me. It also steadied me. A forbidden object could be seized. A dangerous object could be locked away. A teaching object, even under supervision, implied a future in which my mind was more than a breach to be contained.\par

I opened the cover. The first page showed a drawing of a lens. A lens drawn as if anyone might study it, copy it, argue with it, improve it, misunderstand it in public. The obscenity of that freedom was almost beautiful.\par

Va-Elise watched me read the first diagram. ``You are protected here,'' she said.\par

I kept my eyes on the page. ``Not imprisoned?''\par

``Do not force comfort to lie.''\par

Jean-Charles said, ``Protected with locks.''\par

``That is imprisonment.''\par

``Sometimes,'' Va-Elise said. ``Sometimes it is architecture admitting what danger has already made true.''\par

I turned one page. The next drawing showed light entering water. My hand stopped. Va-Elise saw that too. ``We will begin there tomorrow.''\par

Tomorrow. The word entered the room quietly, more dangerous than kindness. Death had no tomorrow. Exception, apparently, did.\par

\ornament

That night I lay in a bed too soft and listened to the house breathe. No hidden pipes. No stone heartbeat. No distant gate chain. A heater clicked. A car passed below the road. Somewhere in the walls, an electrical system hummed with a steadiness I distrusted. The window showed lake darkness and scattered lights along the far shore. They trembled in the water like signals from a city that had no idea it was being watched.\par

I had been Kailan. Then Sa-Kailan. Then Va-Kailan. Now Sa-Kylian. A Candidate again, though Light had no part in it. Candidate of Horizon, perhaps. Candidate of survival. Candidate of proving Shield right to leave me breathing.\par

I touched the place on my wrist where the white Veil knot had been covered. Horizon had wrapped it beneath an ordinary sleeve. That was worse. A removed sign could be mourned. A hidden sign continued to accuse the skin.\par

I thought of my mother receiving a form she could survive. My father hearing what the form withheld. Ka-Raedin managing the office consequence. Va-Sheva returning through desert and stone to a Mountain that would continue without me. Ka-Sarah wearing a cross beneath books and maps. Jean-Charles smiling as if the world's dangers had become personal acquaintances. Va-Elise turning exile into instructions.\par

On the table beside the bed lay the optics manual. Private writing materials stayed forbidden. The lake stayed forbidden. The curtains stayed open until I asked, because Va-Elise had said, ``First learn what windows do.''\par

But the window was there. I sat up and looked through it. Water opened under sky without shame. Houses stood along the far shore with their lights exposed. No Mountain shadow covered them. No Faith attendant named the hour. No Shield warden told the lake where to end.\par

It should have looked like freedom. It looked like exposure. Still, it was a window.\par

Somewhere in that strange geography of outside power there was already a man named Paul Morgan. The name meant nothing to me then. The last custody of my life would arrive as an old man with a recorder and too many questions, in a country of open water. I knew only that the Mountain reached past the Mountain. Its shadow had crossed the desert, entered a dead well, passed into vehicles, aircraft, papers, houses, lakes, churches I was forbidden to enter, and names reshaped for foreign mouths.\par

The horizon is supposed to be where sight ends. For me it was where another custody began. The cage had widened. It stayed shut.\par

But for the first time, light entered from outside.\par

\ornament

The reels have nearly run their length. Four days of him have wound from the left spool onto the right, and what began as bare tape is a full one now --- the boy, the naming, the Pendulum and its nearly three days, the fire in the north, the drowned eye in the west, the Beast under the floor, the round Earth and the hand that reads it, and last of all the road that carried the whole of it out of the stone and set it down, freely, in this house. He has given me the working of it rather than a dot on a map --- the method rather than the answer, which is the harder thing to counterfeit and the harder thing to hold. I came into this house to prove there was no such man. I have four days of him instead, and no honest column left in which to keep the doubt.\par

What escapes him --- what I have only this evening let myself see --- is where the road he has just walked comes out. He left the Mountain and went toward the one shape in all the world he trusted without proof: Horizon, his own side, the friend of the stone. He reached it. He is reaching it still, in this chair, four days deep, handing the last of what he carries to a house whose name is written across the top of my own contract as well as across his hopes. He believes he has come home, and I, who was sent to doubt every word he speaks, find that is the one word I can no longer doubt. What I have begun to doubt instead is whether the home was worth the walking. No man crosses a desert on foot toward people he has never met and arrives, by pure fortune, among friends; he arrives among people who arranged to be found. Four days ago the sentence would have escaped me. Tonight it refuses to leave, and it is my own employers it is about.\par

``That is the end of it,'' he says. ``There is no more Mountain to give you.''\par

``No,'' I say. ``There is not.'' I reach across and stop the reels, and the small sound of it is the loudest thing in the room. ``You have given a good account. All of it. You may rest now.''\par

He believes he has been thanked, and he has been. He looks lighter than he has in four days --- a man who has carried a great weight to the one place he was certain it would be safe, and set it down. I let him look it; the look is honest, whatever else in this room lies. But the tape leaves my hands. It will be copied and carried past the lake to people I have met once and failed to understand, for reasons no one has troubled to give me. And the ease on his face is the ease of a man who has handed his whole life to my house on the strength of a trust I am, tonight, no longer sure I share. I have spent four days proving to myself that everything he says is true. On the last of them the only figure that refuses to come right is the one I have carried since the trial and still lack the courage to total: what a careful, patient, unspeakably rich house wants with a drowned people and a boy who can read light --- and why, for the first time, I find the question frightens me on his account rather than my own.\par

\ornament
